\section{Scope, Limitations and Assumptions}\label{sec:scope-limitations-assumptions}

\subsection*{Scope}

This dissertation is an individual study of a single component, the Intent
Recognition Engine (Component 01 of the NeuroShell framework, project reference
R26-CS-018), and does not implement the downstream planner, executor or
vulnerability-analysis components. Within that boundary, the scope covers:
(1) a guarded NLU pipeline that accepts free-text security requests and emits
validated intent contracts; (2) the deterministic validation and security
countermeasures that constrain the model's output; (3) an offline, CPU-only,
single-device deployment in which nothing is retrieved at inference time and no
request leaves the device; and (4) a closed intent taxonomy and an
operator-configured engagement scope that bounds permitted targets and
techniques. Evaluation is conducted on a curated, annotated data set of prompts
in a lab harness, not against live production traffic.

\subsection*{Limitations}

The principal limitation is evidential: the results reported in this
dissertation rest on a bounded, committed experiment record of
single-run point estimates, without confidence intervals or seed sweeps, and
therefore carry residual uncertainty that is discussed in
Chapter~\ref{ch:discussion}. Only two offline language-model backends are
compared under the guardrail stack, and the evaluation data set is
author-authored rather than drawn from observed operational traffic.
Measurements are specific to the evaluation hardware and CPU-only inference
runtime, so latency and throughput figures are indicative rather than
generalisable. The framework-level benefit of the IRE is validated at the
component level; end-to-end orchestration with the downstream components is
outside the measured evidence presented here.

\subsection*{Assumptions}

The following assumptions underpin the design and evaluation:
\begin{itemize}
\item The operator configures the engagement scope and authorisation policy
correctly; the scope guard can only enforce the boundary it is given.
\item The deployment device and the downloaded model weights are trusted;
physical compromise of the device and backdoors embedded in the model weights
are out of scope.
\item Downstream components consume the intent contract as specified and
honour the contract rule that actions derived for rejected intents must not be
executed.
\item Because the pipeline performs no retrieval and renders no free-form
output, indirect prompt injection via retrieved content and tool injection via
unfiltered rendered text are structurally out of scope, narrowing the attack
surface to the classification of the single user utterance.
\item The annotated prompts in the evaluation data set are representative of
the phrasing an operator would use with the system.
\end{itemize}